\chapter{Sesión 5: Kafka productor, consumidor y notificaciones}
\label{ch:sesion5}

\section{Objetivo de la sesión}

Convertir la publicación de avisos en un flujo de eventos distribuido:

\begin{itemize}[leftmargin=*]
  \item Publicar un evento \texttt{PostCreatedEvent} en Kafka cada vez que se crea un post.
  \item Consumir ese evento y notificar a los suscriptores del canal (excepto al autor).
  \item Mantener contadores de publicaciones no leídas por usuario en Redis.
  \item Exponer endpoints para consultar y limpiar esos contadores.
\end{itemize}

\section{Contexto}

La Sesión~\ref{ch:sesion4} dejó los endpoints de canales y publicaciones funcionando, pero \textbf{las notificaciones eran síncronas y locales}: el post se guardaba en PostgreSQL y la petición HTTP terminaba sin avisar a nadie. Este diseño contradice el modelo pub-sub descrito en el marco teórico (secciones 2.2 y 2.3.1).

Se necesitaba \textbf{desacoplar} la publicación de la notificación: el emisor del post no debe saber quién lo va a leer ni cuándo.

\section{Desarrollo}

\subsection{Estructura creada}

\begin{lstlisting}[language=bash]
src/main/java/mx/uam/notiuam/
|- config/
|  `- KafkaConfig.java                  (definicion del topic)
|- event/
|  |- PostCreatedEvent.java             (payload)
|  |- PostEventPublisher.java           (productor)
|  `- PostNotificationConsumer.java     (consumidor)
|- service/
|  `- UnreadCounterService.java         (contadores en Redis)
`- controller/
   `- MeController.java                 (endpoints de no leidos)
\end{lstlisting}

\subsection{Decisiones técnicas}

\begin{itemize}[leftmargin=*]
  \item \textbf{Topic con 3 particiones.} Permite paralelismo horizontal si más adelante se despliegan varias instancias del consumidor.
  \item \textbf{Clave de mensaje = \texttt{channelId}.} Todos los eventos de un mismo canal caen en la misma partición y se procesan en orden cronológico.
  \item \textbf{Consumer group dedicado} (\texttt{notiuam-notifications}). Permite escalar el consumidor sin perder eventos.
  \item \textbf{Contadores en Redis como hash.} Cada usuario tiene un hash \texttt{notiuam:unread:<userId>} donde cada campo es un \texttt{channelId} y el valor es el contador.
  \item \textbf{Auto-creación de topic.} El bean \texttt{NewTopic} hace que Spring Boot cree el topic al arrancar si no existe, sin depender de scripts externos.
  \item \textbf{No notificar al autor.} Aunque el autor sea suscriptor de su propio canal, no recibe notificación de sus propios posts.
\end{itemize}

\subsection{Componentes clave}

\subsubsection{\texttt{KafkaConfig}}

Define el topic como bean:

\begin{lstlisting}[language=java]
@Configuration
public class KafkaConfig {
    public static final String POSTS_TOPIC = "notiuam.posts";

    @Bean
    public NewTopic postsTopic() {
        return TopicBuilder.name(POSTS_TOPIC)
                .partitions(3)
                .replicas(1)
                .build();
    }
}
\end{lstlisting}

\subsubsection{\texttt{PostCreatedEvent}}

Record inmutable con todo lo necesario para notificar: id del post, canal, autor y contenido.

\begin{lstlisting}[language=java]
public record PostCreatedEvent(
        UUID postId,
        UUID channelId,
        String channelName,
        UUID authorId,
        String authorDisplayName,
        String content,
        Instant createdAt
) {}
\end{lstlisting}

\subsubsection{\texttt{PostEventPublisher}}

Envía el evento de forma asíncrona con \texttt{KafkaTemplate}. Usa \texttt{whenComplete} para loguear el resultado o el error sin bloquear la transacción.

\begin{lstlisting}[language=java]
public void publishPostCreated(PostCreatedEvent event) {
    kafkaTemplate.send(KafkaConfig.POSTS_TOPIC,
                       event.channelId().toString(),
                       event)
            .whenComplete((result, ex) -> {
                if (ex != null) {
                    log.error("Error publicando evento: {}", ex.getMessage());
                } else {
                    log.info("Evento publicado: post={} partition={} offset={}",
                        event.postId(),
                        result.getRecordMetadata().partition(),
                        result.getRecordMetadata().offset());
                }
            });
}
\end{lstlisting}

\subsubsection{\texttt{UnreadCounterService}}

Usa \texttt{StringRedisTemplate} para mantener contadores por usuario y canal.

\begin{lstlisting}[language=java]
public void increment(UUID userId, UUID channelId) {
    redis.opsForHash().increment(KEY_PREFIX + userId, channelId.toString(), 1);
}

public Map<String, Long> getAllCounts(UUID userId) {
    Map<Object, Object> raw = redis.opsForHash().entries(KEY_PREFIX + userId);
    return raw.entrySet().stream().collect(Collectors.toMap(
            e -> e.getKey().toString(),
            e -> Long.parseLong(e.getValue().toString())
    ));
}
\end{lstlisting}

\subsubsection{\texttt{PostNotificationConsumer}}

El listener de Kafka. Por cada evento:

\begin{enumerate}[leftmargin=*]
  \item Busca los suscriptores del canal.
  \item Excluye al autor.
  \item Incrementa el contador de no leídos de cada suscriptor en Redis.
  \item Simula el envío de push (en la Sesión 6/7 se reemplazará con FCM real).
\end{enumerate}

\begin{lstlisting}[language=java]
@KafkaListener(
        topics = KafkaConfig.POSTS_TOPIC,
        groupId = "notiuam-notifications"
)
public void onPostCreated(PostCreatedEvent event) {
    List<Subscription> subs = subscriptionRepository
            .findByChannelId(event.channelId());

    for (Subscription sub : subs) {
        if (sub.getUserId().equals(event.authorId())) continue;
        unreadCounters.increment(sub.getUserId(), event.channelId());
        log.info("[PUSH-SIMULADO] -> user={} | {}", sub.getUserId(), event.channelName());
    }
}
\end{lstlisting}

\subsubsection{Modificación a \texttt{PostService}}

\begin{aviso}[Cambio a un archivo previo]
El \texttt{PostService} de la Sesión~\ref{ch:sesion4} se modificó para publicar el evento después de guardar el post en PostgreSQL.

Se inyectó \texttt{PostEventPublisher} en el constructor y se añadió la llamada al publisher antes del \texttt{return}:

\begin{lstlisting}[language=java]
p = postRepository.save(p);

postEventPublisher.publishPostCreated(new PostCreatedEvent(
        p.getId(), c.getId(), c.getName(),
        author.getId(), author.getDisplayName(),
        p.getContent(), p.getCreatedAt()
));

return toResponse(p);
\end{lstlisting}

Con esto, el post se persiste primero y solo después se emite el evento. Si Kafka falla, el post sigue existiendo y el error solo se registra en los logs.
\end{aviso}

\subsubsection{\texttt{MeController}}

Expone los endpoints para consultar y limpiar los contadores:

\begin{lstlisting}[language=java]
@GetMapping("/unread")            // -> Map<channelId, count>
@DeleteMapping("/unread/{id}")    // -> 204 No Content
\end{lstlisting}

\subsection{Cambios a archivos previos}

\begin{itemize}[leftmargin=*]
  \item \texttt{PostService.java}: se añadió el campo \texttt{postEventPublisher} y una llamada a \texttt{publishPostCreated} dentro de \texttt{create()}.
  \item \texttt{application.yml}: ya tenía la configuración de Kafka desde la Sesión~\ref{ch:sesion2}; no se tocó.
  \item \texttt{GlobalExceptionHandler}: sin cambios.
\end{itemize}

\subsection{Problemas encontrados y solución}

\subsubsection{Método \texttt{create} duplicado en \texttt{PostService}}

Al editar el archivo, el método nuevo se pegó debajo del viejo en lugar de reemplazarlo. La compilación falló con:

\begin{lstlisting}[language=bash]
method create(UUID, PostCreateRequest) is already defined
\end{lstlisting}

\textbf{Solución:} eliminar el método viejo, dejar solo el nuevo que incluye la publicación a Kafka. Este error arrastró decenas de falsos positivos de Lombok (\texttt{cannot find symbol: method getId()}), que desaparecieron al corregir la duplicación.

\subsubsection{\texttt{MeController} faltante}

El endpoint \texttt{/api/me/unread} devolvía \texttt{HTTP 403} sin cuerpo. La causa: el archivo \texttt{MeController.java} no había sido creado, por lo que Spring Security interceptaba la ruta con \texttt{anyRequest().authenticated()} y la rechazaba.

\textbf{Solución:} crear el archivo, reiniciar el backend.

\section{Verificación}

\subsection{Preparar dos usuarios}

\begin{lstlisting}[language=bash]
# Usuario A (autor)
TOKEN_A=$(curl -s -X POST http://localhost:8081/api/auth/login \
  -H "Content-Type: application/json" \
  -d '{"email":"aaron@uam.mx","password":"password123"}' | jq -r .accessToken)

# Usuario B (suscriptor)
TOKEN_B=$(curl -s -X POST http://localhost:8081/api/auth/login \
  -H "Content-Type: application/json" \
  -d '{"email":"betzabet@uam.mx","password":"password123"}' | jq -r .accessToken)
\end{lstlisting}

\subsection{Publicar y observar los logs}

Crear canal, suscribir a ambos usuarios, publicar con el usuario A:

\begin{lstlisting}[language=bash]
CHANNEL_ID=$(curl -s -X POST http://localhost:8081/api/channels \
  -H "Content-Type: application/json" \
  -H "Authorization: Bearer $TOKEN_A" \
  -d '{"name":"Avisos de Ingenieria","category":"ADMINISTRATIVE","publicChannel":true}' \
  | jq -r .id)

curl -s -X POST "http://localhost:8081/api/channels/$CHANNEL_ID/subscribe" \
  -H "Authorization: Bearer $TOKEN_A" > /dev/null

curl -s -X POST "http://localhost:8081/api/channels/$CHANNEL_ID/subscribe" \
  -H "Authorization: Bearer $TOKEN_B" > /dev/null

curl -s -X POST "http://localhost:8081/api/channels/$CHANNEL_ID/posts" \
  -H "Content-Type: application/json" \
  -H "Authorization: Bearer $TOKEN_A" \
  -d '{"content":"Prueba de Kafka"}' | jq
\end{lstlisting}

\begin{exito}[Logs obtenidos]
\begin{lstlisting}[language=bash]
Evento publicado: post=ed496ba0... canal=f4c3bbdd... partition=1 offset=0
Kafka event recibido: post=ed496ba0... canal=Avisos de Ingenieria
[PUSH-SIMULADO] -> user=b8cc6002... | Avisos de Ingenieria | "Prueba de Kafka"
Notificaciones procesadas: 1 de 2 suscriptores
\end{lstlisting}

El \texttt{1 de 2} confirma que el autor (Aaron) no se notifica a sí mismo, y que el otro suscriptor (Betzabet) sí recibe el aviso.
\end{exito}

\subsection{Consultar y limpiar contadores}

\begin{lstlisting}[language=bash]
# Betzabet consulta sus no leidos
curl -s http://localhost:8081/api/me/unread \
  -H "Authorization: Bearer $TOKEN_B" | jq
# -> { "f4c3bbdd-...": 1 }

# Aaron no deberia tener nada
curl -s http://localhost:8081/api/me/unread \
  -H "Authorization: Bearer $TOKEN_A" | jq
# -> {}

# Betzabet marca como leido
curl -s -X DELETE "http://localhost:8081/api/me/unread/$CHANNEL_ID" \
  -H "Authorization: Bearer $TOKEN_B" | jq

# Ya no debe aparecer
curl -s http://localhost:8081/api/me/unread \
  -H "Authorization: Bearer $TOKEN_B" | jq
# -> {}
\end{lstlisting}

\subsection{Verificación en Redis}

\begin{lstlisting}[language=bash]
docker exec -it notiuam-redis redis-cli

# Dentro del CLI
KEYS notiuam:unread:*
HGETALL notiuam:unread:<userId>
exit
\end{lstlisting}

\subsection{Verificación en Kafka UI}

En \url{http://localhost:8080}:

\begin{itemize}[leftmargin=*]
  \item \textbf{Topics} $\rightarrow$ \texttt{notiuam.posts} $\rightarrow$ 3 particiones con mensajes.
  \item \textbf{Consumer Groups} $\rightarrow$ \texttt{notiuam-notifications} $\rightarrow$ offsets por partición con \texttt{LAG = 0}.
\end{itemize}

\section{Resultado}

\begin{exito}[Estado al cierre]
\begin{itemize}[leftmargin=*]
  \item Publicación de posts $\rightarrow$ evento en Kafka $\rightarrow$ consumer $\rightarrow$ contador en Redis, todo automático.
  \item Arquitectura pub-sub real (desacoplamiento entre productor y consumidor).
  \item Persistencia doble: PostgreSQL para los posts, Redis para los contadores.
  \item Endpoints \texttt{/api/me/unread} funcionales.
  \item Base lista para integrar notificaciones push reales con FCM en la Sesión 6.
\end{itemize}
\end{exito}

\textbf{Pendiente para la siguiente sesión:} cliente móvil en Flutter (login, exploración de canales, suscripción, publicación, contadores, notificaciones push reales).
